\chapter{Минимальная finite-module архитектура endpoint-расширения}
\label{ch:v9-endpoint-finite-module}

\section{Два обязательных матричных блока}

Raw-origin no-go выделил три новые system-линии. На нейтральной паре
$H_{\rm n}=\mathbb Cs_0\oplus\mathbb Ca_0$ требуется матричная единица
$|s_0\rangle\langle a_0|$ и её сопряжение. Минимальное ассоциативное
замыкание поэтому равно
\begin{equation}
 \mathcal A_{\rm n}=M_2(\mathbb C),
 \qquad \dim_{\mathbb C}\mathcal A_{\rm n}=4.
 \label{eq:v9-finite-module-neutral-algebra}
\end{equation}

Новая charged-линия вместе с двумя старыми положительно градуированными
координатами образует
\begin{equation}
 H_{\rm t}=\operatorname{span}_{\mathbb C}
 \{e_R^{(t,0)},e_R^{(t,1)},e_R^{(t,2)}\}\simeq\mathbb C^3.
 \label{eq:v9-finite-module-triplet}
\end{equation}
Family-ковариантное замыкание coordinate projectors и стандартного
$SO(3)$-действия есть полный блок
\begin{equation}
 \mathcal A_{\rm t}=M_3(\mathbb C),
 \qquad \dim_{\mathbb C}\mathcal A_{\rm t}=9.
 \label{eq:v9-finite-module-triplet-algebra}
\end{equation}
Таким образом, минимальная новая endpoint-алгебра имеет вид
\begin{equation}
 \mathcal A_{\min}=M_2(\mathbb C)\oplus M_3(\mathbb C),
 \qquad \dim_{\mathbb C}\mathcal A_{\min}=13,
 \qquad \dim Z(\mathcal A_{\min})=2.
 \label{eq:v9-finite-module-minimal-algebra}
\end{equation}

\section{Точное вложение старого carrier}

Выделим в $H_{21}$ две старые положительные charged-линии и обозначим их
ортогональное дополнение через $H_{\rm rest}$, $\dim H_{\rm rest}=19$.
Тогда один и тот же $H_{24}$ записывается как
\begin{equation}
 H_{24}=H_{\rm rest}\oplus H_{\rm n}\oplus H_{\rm t},
 \qquad 19+2+3=24.
 \label{eq:v9-finite-module-h24-splitting}
\end{equation}
Старое пространство является углом
\begin{equation}
 H_{21}=H_{\rm rest}\oplus
 \operatorname{span}_{\mathbb C}\{e_R^{(t,1)},e_R^{(t,2)}\}
 \subset H_{24}.
 \label{eq:v9-finite-module-old-corner}
\end{equation}
На новом пятикоординатном блоке старая Hermitian endpoint-алгебра содержала
только два charged coordinate projectors. Полные блоки имеют $4+9=13$
Hermitian directions, поэтому точный прирост равен
\begin{equation}
 \Delta N_{\rm Herm}=13-2=11.
 \label{eq:v9-finite-module-hermitian-increment}
\end{equation}

\section{Gauge, grading и family compatibility}

На нейтральном блоке
\begin{equation}
 \Gamma_{\rm n}=\operatorname{diag}(1,-1),\qquad
 Y_{\rm n}=0,\qquad
 X=E_{12}+E_{21},\quad Y=-i(E_{12}-E_{21})
 \label{eq:v9-finite-module-neutral-data}
\end{equation}
и точно
\begin{equation}
 \{\Gamma_{\rm n},X\}=\{\Gamma_{\rm n},Y\}=0,
 \qquad [Y_{\rm n},X]=[Y_{\rm n},Y]=0.
 \label{eq:v9-finite-module-neutral-compatibility}
\end{equation}

На triplet-блоке
\begin{equation}
 \Gamma_{\rm t}=I_3,\qquad Y_{\rm t}=-I_3,
 \qquad [J_a,\Gamma_{\rm t}]=[J_a,Y_{\rm t}]=0
 \quad(a=1,2,3).
 \label{eq:v9-finite-module-triplet-compatibility}
\end{equation}
Следовательно, полный $M_3$-блок совместим одновременно с family,
grading и hypercharge.

\section{Real-замыкание}

Charged endpoint требует сопряжённого представления. Оно добавляется не
как новый свободный selector, а функториально:
\begin{equation}
 \widehat H_{24}=H_{24}\oplus\overline{H}_{24},\qquad
 \pi_R(a)=\pi(a)\oplus\overline{\pi(a)},\qquad
 J(v,w)=(\overline w,\overline v).
 \label{eq:v9-finite-module-real-double}
\end{equation}
Отсюда
\begin{equation}
 J^2=1,\qquad J\pi_R(a)J^{-1}=\pi_R(a),
 \qquad
 \dim_{\mathbb C}(\mathcal E_{\min}\oplus J\mathcal E_{\min})=6.
 \label{eq:v9-finite-module-real-identities}
\end{equation}
Независимых новых system-данных при этом остаётся три complex lines:
вторая половина Real-дубля определяется первой.

\section{Архитектурный итог}

Десять точных проверок дают
\begin{equation}
 N_{\rm finite\ module\ architecture}=10/10.
 \label{eq:v9-finite-module-architecture-ledger}
\end{equation}
Но новый summand $\mathcal A_{\min}$ пока введён как минимальная
архитектура, а не получен стационарным условием прежнего действия:
\begin{equation}
 N_{\rm finite\ module\ parent\ origin}=0/1,\qquad
 N_{\rm raw\ physical\ slot\ closure}=0/4.
 \label{eq:v9-finite-module-origin-ledger}
\end{equation}

\begin{theorem}[Минимальная finite-module архитектура]
Для одновременного neutral linking и family-triplet completion необходимо и
достаточно присоединить endpoint-алгебру
$M_2(\mathbb C)\oplus M_3(\mathbb C)$ на пятикоординатном активном блоке.
Она требует ровно три новые complex system-линии, добавляет одиннадцать
Hermitian directions, сохраняет старый $H_{21}$ как угол и допускает
канонический Real-дубль. Архитектура минимальна, но её parent-origin не
получен.
\end{theorem}

\begin{nogocriterion}[Минимальная алгебра не есть динамическое рождение]
Факт, что $M_2\oplus M_3$ является наименьшим замыканием требуемых matrix
units, не доказывает, что этот summand возникает из вакуума старого action.
До вывода его projector или rank-changing order parameter новый module
остаётся явным typed input.
\end{nogocriterion}

\begin{protocol}\begin{enumerate}
\item neutral algebra: \textbf{$M_2(\mathbb C)$};
\item family endpoint algebra: \textbf{$M_3(\mathbb C)$};
\item complex dimension новой алгебры: \textbf{$13$};
\item centre dimension: \textbf{$2$};
\item новых independent complex states: \textbf{$3$};
\item Real-замкнутый прирост: \textbf{$6$ complex coordinates};
\item новых Hermitian directions: \textbf{$11$};
\item old $H_{21}$ сохраняется как угол: \textbf{да};
\item architecture ledger: \textbf{$10/10$};
\item parent-origin нового module: \textbf{$0/1$};
\item следующий гейт: \textbf{parent-action origin finite endpoint-module}.
\end{enumerate}\end{protocol}

Машинный сертификат сохранён в
\path{s2t_v9_endpoint_extension_minimal_finite_module_architecture_gate_results.json}.